\SourceLeaf{377}%
die Signale und erschweren die Beobachtungen selbst erheblich. Diese Umstände veranlassten die Vermesser, verhältnismäßig viele recht kleine Dreiecke zu bilden. Es sind sechzehn, und die kürzeste Verbindungsseite ist nur 18671 Meter lang. Wir dürfen jedoch nicht unerwähnt lassen, dass Heliotrope damals in Italien noch unbekannt waren. Sie machen alle Signale ganz entbehrlich und erlauben es, die Winkel in den allergrößten Dreiecken ebenso leicht und genau zu messen wie im kleinsten.

Die Winkel wurden mit achtzölligen Theodoliten von \textsc{Reichenbach} gemessen. Die piemontesischen und österreichischen Offiziere arbeiteten getrennt und bestimmten jeweils alle Winkel des Netzes unabhängig voneinander. Man erhielt also für jeden einzelnen Winkel zwei Bestimmungen und nahm aus ihnen, unter Berücksichtigung der Zahl der jeweils eingegangenen Messreihen, den Mittelwert als endgültigen Wert an. Die Ergebnisse der piemontesischen Offiziere beruhen meist auf sechs Reihen mit je 10 Wiederholungen, die der österreichischen größtenteils auf zwei, einige auf drei oder vier Reihen. Alle Messungen sind sehr ausführlich abgedruckt. Wer die einzelnen Beobachtungen angestellt hat, ist allerdings nicht angegeben.

Bei einer so umfangreichen Vermessung ist die erreichte Genauigkeit der Winkelbestimmungen von großem Interesse. Die Winkelsummen der einzelnen Dreiecke bieten dafür einen Anhaltspunkt, der unter Umständen allerdings etwas trügerisch sein kann. Darf man die vorliegenden Messungen danach beurteilen, so ist ihre Genauigkeit in der Tat bewundernswert. Der größte Fehler der Winkelsumme bei den 16 Dreiecken beträgt nur \(1''16\). Der mittlere Fehler beträgt \(0''70\), und der mittlere Fehler der einzelnen Winkel wäre demnach nur \(0''40\). Möglichkeiten zur Prüfung durch Diagonalrichtungen oder durch Bildung von Polygonen fehlen ganz. Der Vergleich der beiden Bestimmungen jedes der 48 Winkel deutet aber ganz entschieden auf eine bei Weitem größere Ungenauigkeit der Ergebnisse hin. Wir finden hier 13 Fälle, in denen der Unterschied über \(3''\) liegt, darunter 5, in denen er über \(5''\) steigt. In einem Fall, gleich im ersten Dreieck, weicht die auf 80 Wiederholungen gestützte Bestimmung der piemontesischen Offiziere von der auf 48 Wiederholungen beruhenden Bestimmung der österreichischen Offiziere sogar um \(9''2\) ab. Bei so großen Unterschieden drängt sich die Vermutung auf, dass eine richtige Beurteilung der tatsächlichen Genauigkeit der Messungen auch von weiteren Umständen abhängt, über die das Werk uns nichts mitteilt.

